\section{Discussion and Limitations}\label{sec:discuss}

\paragraph*{Operating points}
M-CTX deliberately exposes a small menu rather than a single index. For the
OSM stage, the STR-tree is the right default---exact, $23$\,$\mu$s,
$349$\,KB, zero native dependencies; a deployment that already runs a
production R$^{*}$-tree may prefer it for raw query throughput; and BR-LZ
is the choice where a \emph{recall guarantee}, the smallest learned
footprint, and the fastest build matter and per-query cost is amortised
over the corpus. For the SDF stage the recommendation is unambiguous:
\texttt{uint8} at $32\!\times\!32$, no clipping ($64\times$ compression at
$0.04$\,m ADE). For neighbours, the B$^{x}$-tree is mandatory under
streaming and the KD-tree is marginally cheaper only in a fully static
batch. This menu lets a deployment trade build, latency, footprint, and
formal guarantees explicitly, which is what a real ingest pipeline needs.

\paragraph*{Limitations}
M-CTX targets read-mostly OSM; workloads with frequent \emph{feature}
updates would need an insertion-friendly index family, which we do not
address. The B$^{x}$-tree assumes a snapshot cadence; at sub-second update
rates the phase encoding would require redesign to avoid key-space churn.
Our reported numbers are single-node in-memory; we capture the
partition and routing costs of a cluster with a space-tiled shard
simulator (Appendix~\ref{app:extended}), but inter-node RTT is left for a
multi-node study. The $40$\,M-feature scale-up uses deterministic spatial
replication of the DMA corpus; a real planet-scale OSM extract with a
memory-mapped backing store would re-validate the curve at higher scales.
BR-LZ's vectorised back-end delivers $100$\,$\mu$s $p_{50}$ on DMA and $2.57$\,ms at $N{=}40$\,M synthetic; the asymptotics and correctness are preserved (segment layout unchanged from Theorem~\ref{thm:brlz-recall} and Proposition~\ref{prop:pareto}).  A C/Cython port would further narrow the remaining constant-factor gap to libspatialindex.
